%! AP Statistics — Unit 5, Lesson 5.2 — Lesson Plan
\documentclass[10pt]{article}
\usepackage{apstats-article}
\usepackage{apstats-boxes}
\newcommand{\TallMath}[1]{$\displaystyle #1\rule[-1.4em]{0pt}{3.2em}$}
\newcommand{\UnitNumberName}{Unit 5: Sampling Distributions \quad}
\newcommand{\LessonNumberName}{Lesson 5.2: The Normal Distribution, Revisited}

\begin{document}
\begin{center}
    {\Large\bfseries \CourseName: \SchoolYear \\
    {\small\bfseries \UnitNumberName \LessonNumberName}}
\end{center}

\begin{tcolorbox}[colback=sky, colframe=navy, boxrule=0.9pt, arc=2mm,
    left=2mm, right=2mm, top=1.5mm, bottom=1.5mm]
    \textbf{Primary Objective:} Students will calculate probabilities for a normal
    distribution using z-scores and technology, determine intervals associated with
    given areas, and assess when the normal distribution is an appropriate model
    (Big Idea: VAR).
\end{tcolorbox}

\renewcommand{\arraystretch}{1.6}
\begin{skillbox}[Priority Ideas \& Skills]{goldbox}
\begin{minipage}[t]{0.48\linewidth}
    \textbf{AP Skill 3 --- Using Probability and Simulation}
    \begin{itemize}
        \item Calculate probabilities for normal distributions using z-scores or
              technology (Skill 3.A; VAR-6.A).
        \item Determine boundary values for a given area using z-scores or a standard
              normal table (Skill 3.A; VAR-6.B).
        \item Assess whether the normal distribution is appropriate to model an unknown
              distribution (Skill 3.C; VAR-6.C).
    \end{itemize}
\end{minipage}
\hfill
\begin{minipage}[t]{0.48\linewidth}
    \textbf{Key Understandings}
    \begin{itemize}
        \item A continuous random variable can take any value in a domain; area under
              the normal curve equals probability (VAR-6.A.1--3).
        \item Boundary values of a normal interval are found with z-scores or
              technology (VAR-6.B.1--2).
        \item Normal distributions can approximate other symmetric, bell-shaped
              distributions (VAR-6.C.1).
        \item This lesson is preparation for the sampling distributions in Lessons
              5.5--5.8, which are all approximately normal.
    \end{itemize}
\end{minipage}
\end{skillbox}

\begin{skillbox}[{Vocabulary, Concepts, \& Theorems}]{greenbox}
\renewcommand{\arraystretch}{1.3}
\begin{tabularx}{\linewidth}{@{}lX@{}}
    \textbf{Continuous random variable} & A variable that can take any value in a
    specified domain; every interval has an associated probability.\\
    \textbf{Normal distribution} & A symmetric, bell-shaped probability distribution
    described by mean $\mu$ and standard deviation $\sigma$.\\
    \textbf{z-score} & \TallMath{z = \dfrac{x - \mu}{\sigma}} --- number of standard
    deviations $x$ is from the mean.\\
    \textbf{Standard normal} & Normal distribution with $\mu = 0$, $\sigma = 1$.\\
    \textbf{Area under curve} & Equals the probability that $X$ falls in that interval.\\
\end{tabularx}
\end{skillbox}

\begin{skillbox}[Activate Prior Knowledge \& Spiral Review]{sky}
\begin{minipage}[t]{0.48\linewidth}
    \textbf{Prior Knowledge Check (3--4 min)}
    \begin{itemize}
        \item Review: Shape, center, and spread of a normal distribution (Unit 1).
        \item Recall: How to compute a z-score (Unit 1, Lesson 1.9).
        \item Prompt: \textit{``SAT scores are approximately normal with $\mu = 1060$
              and $\sigma = 195$. What is the z-score for a score of 1255? What does
              that tell us?''}
    \end{itemize}
\end{minipage}
\hfill
\begin{minipage}[t]{0.48\linewidth}
    \textbf{Connections Forward}
    \begin{itemize}
        \item Normal probabilities for a population $\to$ today
        \item Sampling distributions are approximately normal $\to$ Lessons 5.5--5.8
        \item Use normal model for inference on proportions $\to$ Unit 6
        \item Use normal model for inference on means $\to$ Unit 7
    \end{itemize}
\end{minipage}
\end{skillbox}

\begin{skillbox}[Hook \& Lesson]{sky}
\begin{minipage}[t]{0.48\linewidth}
    \textbf{Hook (5 min)}\\
    Display a histogram of adult heights in the U.S. (approximately normal,
    $\mu = 69.2$ in, $\sigma = 3.6$ in for males).\\[4pt]
    Ask: ``What fraction of adult males are taller than 6 feet (72 in)?'' Let students
    estimate from the histogram.\\[4pt]
    Reveal: We can compute this exactly using the normal model --- and the answer is
    about 19\%. Today's lesson shows how.
\end{minipage}
\hfill
\begin{minipage}[t]{0.48\linewidth}
    \textbf{Lesson Flow (18--20 min)}
    \begin{enumerate}
        \item \textbf{Review:} $X \sim N(\mu, \sigma)$ notation; bell curve symmetry;
              the 68--95--99.7 rule.
        \item \textbf{Forward probability} (area to the left/right of a value):
              compute $z$, then use table or technology.
        \item \textbf{Middle area} (area between two values): subtract two left-tail
              probabilities.
        \item \textbf{Inverse normal} (finding $x$ given area): work backward using
              z-score formula or \texttt{invNorm}.
        \item \textbf{Two-sided area:} identifying most extreme $p\%$ requires
              splitting into two tails (VAR-6.B.2d).
        \item Assess appropriateness: use normal only when distribution is
              approximately symmetric and bell-shaped (VAR-6.C.1).
    \end{enumerate}
\end{minipage}
\end{skillbox}

\begin{skillbox}[Explicit Instruction \& Active Monitoring]{sky}
\begin{minipage}[t]{0.48\linewidth}
    \textbf{Worked Example: Forward Probability}\\
    $X \sim N(69.2, 3.6)$. Find $P(X > 72)$.
    \begin{enumerate}
        \item Compute $z = \frac{72 - 69.2}{3.6} \approx 0.78$.
        \item $P(Z > 0.78) = 1 - 0.7823 = 0.2177$.
        \item Interpret: About 22\% of adult males are taller than 72 in.
    \end{enumerate}
    \textbf{Worked Example: Inverse Normal}\\
    What height is the 90th percentile?
    \begin{enumerate}
        \item Find $z$ with $P(Z < z) = 0.90$ $\Rightarrow$ $z \approx 1.28$.
        \item $x = \mu + z\sigma = 69.2 + 1.28(3.6) \approx 73.8$ in.
    \end{enumerate}
\end{minipage}
\hfill
\begin{minipage}[t]{0.48\linewidth}
    \textbf{Active Monitoring}
    \begin{itemize}
        \item Watch for students who do not sketch the curve and shade --- require a
              sketch and label of the mean before computing.
        \item Cold-call: ``Is this a left-tail or right-tail probability? How do you
              know?''
        \item Check that students subtract from 1 for right-tail areas and divide by 2
              for two-tailed ``most extreme'' problems.
        \item Common error: using the z-score itself as the probability --- redirect to
              the table or \texttt{normalcdf}.
    \end{itemize}
\end{minipage}
\end{skillbox}

\begin{skillbox}[Group Work \& Differentiation]{redbox}
\begin{minipage}[t]{0.48\linewidth}
    \textbf{Partner Activity (10 min)}\\
    Three-context activity (birth weights, standardized test scores, daily rainfall):
    \begin{enumerate}
        \item Calculate a forward probability (shade and label).
        \item Calculate a probability between two values.
        \item Find a value given a percentile (inverse normal).
    \end{enumerate}
\end{minipage}
\hfill
\begin{minipage}[t]{0.48\linewidth}
    \textbf{Differentiation}
    \begin{itemize}
        \item \textit{Support (Tier R)}: Provide a z-table and a formula card; focus on
              steps 1 and 2 only.
        \item \textit{On-grade (Tier A)}: Complete all three problems; write full
              contextual interpretations.
        \item \textit{Extension (Tier E)}: Solve the most-extreme-$p\%$ problem (two
              tails); critique whether the normal model is appropriate for the context.
        \item \textit{ELL}: Label curve diagram vocabulary in students' first language.
    \end{itemize}
\end{minipage}
\end{skillbox}

\begin{skillbox}[Individual Work \& Assessment]{redbox}
\begin{minipage}[t]{0.48\linewidth}
    \textbf{Exit Ticket (5 min)}\\
    SAT scores $\sim N(1060, 195)$.
    \begin{enumerate}
        \item Find the probability that a randomly selected student scored above 1300.
        \item Find the SAT score at the 75th percentile.
        \item Is it appropriate to use the normal model here? Explain.
    \end{enumerate}
\end{minipage}
\hfill
\begin{minipage}[t]{0.48\linewidth}
    \textbf{AP-Style Check}\\
    \textit{$X \sim N(50, 8)$. Which is closest to $P(34 < X < 66)$?}
    \begin{enumerate}[label=(\Alph*)]
        \item 0.68
        \item 0.95
        \item 0.997
        \item 0.50
    \end{enumerate}
    Answer: (B) (two SDs on each side $\approx 0.95$)
\end{minipage}
\end{skillbox}

\begin{skillbox}[Reinforcement \& Extension]{goldbox}
\begin{minipage}[t]{0.48\linewidth}
    \textbf{Homework / Practice}
    \begin{itemize}
        \item Six normal-probability problems covering forward, backward, and between
              scenarios across three contexts (heights, weights, exam scores).
        \item One ``appropriateness'' question: given a histogram, decide whether a
              normal model is appropriate and justify.
    \end{itemize}
\end{minipage}
\hfill
\begin{minipage}[t]{0.48\linewidth}
    \textbf{Extension / Preview}
    \begin{itemize}
        \item \textit{Extension}: Investigate the effect of $\sigma$ on normal
              probability: for fixed $\mu = 100$, compare $P(X > 110)$ when
              $\sigma = 5$, 10, 20. What pattern do you notice?
        \item \textit{Preview}: Lesson 5.3 introduces the Central Limit Theorem ---
              which states that sample means form an approximately normal distribution
              for large enough samples, even when the population isn't normal.
    \end{itemize}
\end{minipage}
\end{skillbox}

\end{document}
